% GENERATED FILE. Edit canonical lessons, then run npm run generate:books.
% canonical-source-hash: fnv1a64:b9c92881368ce5b8
% canonical-lessons: AR-C351-khadaa-trick, AR-C351-ghashsha, AR-C351-ihtala, AR-C351-zawwara, AR-C351-akhfa

\chapter{Tricks and lies}
\label{ch:ar-tricks-and-lies}
\hlchaptermodality{351}

\begin{chapteropening}
\textbf{By the end of this chapter:} \emph{I can say to deceive, to cheat, to swindle, to forge and to conceal.}

Complete the last lesson of chapter 351: I can say to deceive, to cheat, to swindle, to forge and to conceal.
\end{chapteropening}

\section[\texorpdfstring{\ar{خدع} (khadaʿa)}{خدع (khadaʿa)}]{\ar{خدع} (khadaʿa) — to deceive, to trick}

\label{lesson:AR-C351-khadaa-trick}



\begin{quote}
Before the new one: say the Arabic for a leaflet, then the Arabic for a supplement.
\end{quote}

\subsection*{You'll want to know: \ar{خدع}}
\textbf{\ar{خدع}} — \emph{khadaʿa} — \textquotedblleft{}to deceive, to trick\textquotedblright{}.

\begin{quote}
\ar{خدع} \ar{الثعلب} \ar{الغراب} \ar{في} \ar{الحكاية}.
\end{quote}

\begin{quote}
\emph{khadaʿa ath-thaʿlab al-ghurāb fī al-ḥikāya.}
\end{quote}

\begin{quote}
The fox tricked the crow in the tale.
\end{quote}

Present: \ar{يخدع}, \emph{yakhdaʿ}.

A trick is \ar{خدعة}.

\begin{cousinweb}
\textbf{\ar{خدع}} (\emph{khadaʿa}) is built on the native Arabic root \emph{kh-d-ʿ} (Lane), \textquotedblleft{}to hide a thing; to deceive\textquotedblright{}; \emph{khidʿa} is a ruse.
\end{cousinweb}

\begin{grammarlens}[title={What you've built}]
One more word for telling what happened.
\end{grammarlens}

\subsection*{Your turn}
\begin{itemize}
\raggedright
  \item \emph{Say it:} \emph{khadaʿa}
  \item \emph{Say it:} \emph{khadaʿa}, once more
  \item \emph{Say it:} \emph{khadaʿa ath-thaʿlab al-ghurāb fī al-ḥikāya.} — The fox tricked the crow in the tale.
  \item \emph{Recall:} say the Arabic for a leaflet, then the Arabic for a supplement, then say \emph{khadaʿa} again
\end{itemize}

\begin{tcolorbox}[breakable,colback=teal!4,colframe=teal!35!black,title={Before you move on}]
What does \ar{خدع} mean? (\textquotedblleft{}To deceive, to trick\textquotedblright{}.) Say it once more.
\end{tcolorbox}
\section[\texorpdfstring{\ar{غش} (ghashsha)}{غش (ghashsha)}]{\ar{غش} (ghashsha) — to cheat}

\label{lesson:AR-C351-ghashsha}



\begin{quote}
Before the new one: say the Arabic for a supplement, then the Arabic for to deceive.
\end{quote}

\subsection*{You'll want to know: \ar{غش}}
\textbf{\ar{غش}} — \emph{ghashsha} — \textquotedblleft{}to cheat\textquotedblright{}.

\begin{quote}
\ar{غش} \ar{الطالب} \ar{في} \ar{الامتحان}.
\end{quote}

\begin{quote}
\emph{ghashsha aṭ-ṭālib fī al-imtiḥān.}
\end{quote}

\begin{quote}
The student cheated in the exam.
\end{quote}

Present: \ar{يغش}, \emph{yaghushsh}.

Cheating is \ar{غش}.

\begin{cousinweb}
\textbf{\ar{غش}} (\emph{ghashsha}) is built on the native Arabic root \emph{gh-sh-sh} (Lane), \textquotedblleft{}to be insincere towards someone\textquotedblright{}, to show what one does not mean.
\end{cousinweb}

\begin{grammarlens}[title={What you've built}]
One more word for telling what happened.
\end{grammarlens}

\subsection*{Your turn}
\begin{itemize}
\raggedright
  \item \emph{Say it:} \emph{ghashsha}
  \item \emph{Say it:} \emph{ghashsha}, once more
  \item \emph{Say it:} \emph{ghashsha aṭ-ṭālib fī al-imtiḥān.} — The student cheated in the exam.
  \item \emph{Recall:} say the Arabic for a supplement, then the Arabic for to deceive, then say \emph{ghashsha} again
\end{itemize}

\begin{tcolorbox}[breakable,colback=teal!4,colframe=teal!35!black,title={Before you move on}]
What does \ar{غش} mean? (\textquotedblleft{}To cheat\textquotedblright{}.) Say it once more.
\end{tcolorbox}
\section[\texorpdfstring{\ar{احتال} (iḥtāla)}{احتال (iḥtāla)}]{\ar{احتال} (iḥtāla) — to swindle}

\label{lesson:AR-C351-ihtala}



\begin{quote}
Before the new one: say the Arabic for to deceive, then the Arabic for to cheat.
\end{quote}

\subsection*{You'll want to know: \ar{احتال}}
\textbf{\ar{احتال}} — \emph{iḥtāla} — \textquotedblleft{}to swindle\textquotedblright{}.

\begin{quote}
\ar{احتال} \ar{رجل} \ar{على} \ar{السياح} \ar{في} \ar{السوق}.
\end{quote}

\begin{quote}
\emph{iḥtāla rajul ʿalā as-suyyāḥ fī as-sūq.}
\end{quote}

\begin{quote}
A man swindled the tourists in the market.
\end{quote}

Present: \ar{يحتال}, \emph{yaḥtāl}.

It takes \ar{على}. A swindler is \ar{محتال}.

\begin{cousinweb}
\textbf{\ar{احتال}} (\emph{iḥtāla}) is built on the native Arabic root \emph{ḥ-w-l} (Lane), \textquotedblleft{}to change, to shift\textquotedblright{}, where \emph{ḥīla} is a shift or wile. It is Form VIII, \textquotedblleft{}to use a wile\textquotedblright{}.
\end{cousinweb}

\begin{grammarlens}[title={What you've built}]
One more word for telling what happened.
\end{grammarlens}

\subsection*{Your turn}
\begin{itemize}
\raggedright
  \item \emph{Say it:} \emph{iḥtāla}
  \item \emph{Say it:} \emph{iḥtāla}, once more
  \item \emph{Say it:} \emph{iḥtāla rajul ʿalā as-suyyāḥ fī as-sūq.} — A man swindled the tourists in the market.
  \item \emph{Recall:} say the Arabic for to deceive, then the Arabic for to cheat, then say \emph{iḥtāla} again
\end{itemize}

\begin{tcolorbox}[breakable,colback=teal!4,colframe=teal!35!black,title={Before you move on}]
What does \ar{احتال} mean? (\textquotedblleft{}To swindle\textquotedblright{}.) Say it once more.
\end{tcolorbox}
\section[\texorpdfstring{\ar{زور} (zawwara)}{زور (zawwara)}]{\ar{زور} (zawwara) — to forge, to falsify}

\label{lesson:AR-C351-zawwara}



\begin{quote}
Before the new one: say the Arabic for to cheat, then the Arabic for to swindle.
\end{quote}

\subsection*{You'll want to know: \ar{زور}}
\textbf{\ar{زور}} — \emph{zawwara} — \textquotedblleft{}to forge, to falsify\textquotedblright{}.

\begin{quote}
\ar{زور} \ar{الرجل} \ar{جواز} \ar{سفر}.
\end{quote}

\begin{quote}
\emph{zawwara ar-rajul jawāz safar.}
\end{quote}

\begin{quote}
The man forged a passport.
\end{quote}

Present: \ar{يزور}, \emph{yuzawwir}.

Lane files the taught \ar{زار}, to visit, under the same letters z-w-r; the two verbs differ in form and sense.

\begin{cousinweb}
\textbf{\ar{زور}} (\emph{zawwara}) is built on the native Arabic root \emph{z-w-r} (Lane), where \emph{zūr} is falsehood. It is Form II, \textquotedblleft{}to make a false thing look true\textquotedblright{}.
\end{cousinweb}

\begin{grammarlens}[title={What you've built}]
One more word for telling what happened.
\end{grammarlens}

\subsection*{Your turn}
\begin{itemize}
\raggedright
  \item \emph{Say it:} \emph{zawwara}
  \item \emph{Say it:} \emph{zawwara}, once more
  \item \emph{Say it:} \emph{zawwara ar-rajul jawāz safar.} — The man forged a passport.
  \item \emph{Recall:} say the Arabic for to cheat, then the Arabic for to swindle, then say \emph{zawwara} again
\end{itemize}

\begin{tcolorbox}[breakable,colback=teal!4,colframe=teal!35!black,title={Before you move on}]
What does \ar{زور} mean? (\textquotedblleft{}To forge, to falsify\textquotedblright{}.) Say it once more.
\end{tcolorbox}
\section[\texorpdfstring{\ar{أخفى} (akhfā)}{أخفى (akhfā)}]{\ar{أخفى} (akhfā) — to conceal, to hide}

\label{lesson:AR-C351-akhfa}



\begin{quote}
Before the new one: say the Arabic for to swindle, then the Arabic for to forge.
\end{quote}

\subsection*{You'll want to know: \ar{أخفى}}
\textbf{\ar{أخفى}} — \emph{akhfā} — \textquotedblleft{}to conceal, to hide\textquotedblright{}.

\begin{quote}
\ar{أخفت} \ar{البنت} \ar{الهدية} \ar{تحت} \ar{السرير}.
\end{quote}

\begin{quote}
\emph{akhfat al-bint al-hadiyya taḥta as-sarīr.}
\end{quote}

\begin{quote}
The girl hid the present under the bed.
\end{quote}

Present: \ar{يخفي}, \emph{yukhfī}.

\begin{cousinweb}
\textbf{\ar{أخفى}} (\emph{akhfā}) is built on the native Arabic root \emph{kh-f-y} (Lane), \textquotedblleft{}to be hidden\textquotedblright{}, the root of the taught \emph{ikhtafā}, to disappear. It is Form IV, \textquotedblleft{}to hide a thing\textquotedblright{}.
\end{cousinweb}

\begin{grammarlens}[title={What you've built}]
One more word for telling what happened.
\end{grammarlens}

\subsection*{Your turn}
\begin{itemize}
\raggedright
  \item \emph{Say it:} \emph{akhfā}
  \item \emph{Say it:} \emph{akhfā}, once more
  \item \emph{Say it:} \emph{akhfat al-bint al-hadiyya taḥta as-sarīr.} — The girl hid the present under the bed.
  \item \emph{Recall:} say the Arabic for to swindle, then the Arabic for to forge, then say \emph{akhfā} again
\end{itemize}

\begin{tcolorbox}[breakable,colback=teal!4,colframe=teal!35!black,title={Before you move on}]
What does \ar{أخفى} mean? (\textquotedblleft{}To conceal, to hide\textquotedblright{}.) Say it once more.
\end{tcolorbox}
